\documentclass[../../../include/open-logic-section]{subfiles}

\begin{document}

\olfileid[fr]{sth}{card-arithmetic}{fix}
\olsection{Points fixes de la fonction \texorpdfstring{$\aleph$}{ℵ}}

Dans le \olref[spine][]{chap}, nous avons suggéré
que le remplacement demande une justification,
car il impose à la hiérarchie d'atteindre une
hauteur considérable. Maintenant que nous
disposons de quelques résultats d'arithmétique
cardinale, nous pouvons en donner une illustration.

Il est clair que $0 < \aleph_0$, que
$1 < \aleph_1$, que $2 < \aleph_2$\ldots
et, de fait, l'écart de taille devient
\emph{plus grand} à chaque étape.
Il est donc tentant de conjecturer que
$\kappa< \aleph_\kappa$ pour tout ordinal
$\kappa$.

Mais cette conjecture est \emph{fausse}
dans $\ZFC$. Nous pouvons en effet démontrer
qu'il existe des \emph{points fixes de la
fonction $\aleph$}, c'est-à-dire des cardinaux
$\kappa$ tels que $\kappa=\aleph_\kappa$.

\begin{prop}\ollabel{alephfixed}
Il existe un point fixe de la fonction $\aleph$.
\end{prop}

\begin{proof}
Définissons par récurrence :
\begin{align*}
	\kappa_0 &= 0\\
	\kappa_{n+1} &= \aleph_{\kappa_n}\\
	\kappa&= \bigcup_{n < \omega}\kappa_n
\end{align*}
Alors $\kappa$ est un cardinal par la
\olref[cardinals][classing]{unioncardinalscardinal}.
Mais :
\[
	\kappa= \bigcup_{n < \omega} \kappa_{n+1} =
	\bigcup_{n < \omega}\aleph_{\kappa_n} =
	\bigcup_{\alpha < \kappa}\aleph_\alpha
	= \aleph_\kappa.
\]
% Par construction, kappa est le plus petit cardinal dépassant
% chacun des kappa_n ; le même raisonnement s'applique aux alephs.
\end{proof}

Boolos a consacré un article précisément au
point fixe de la fonction $\aleph$ que nous
venons de construire. Après avoir constaté
l'existence de $\kappa$, il écrivait au début
de cet article :
\begin{quote}
	[$\kappa$ est] un nombre \emph{assez grand}
	aux yeux de ceux qui n'ont jamais étudié
	la théorie des ensembles, si grand, me
	semble-t-il, qu'il conduit à mettre en
	question la vérité de toute théorie qui
	affirme qu'il existe au moins $\kappa$
	objets.
	\citep[p.~257]{Boolos2000}
\end{quote}
Il terminait finalement son article
par cette question :
\begin{quote}
	Ne soupçonnons-nous pas que, quelle
	qu'ait été la situation au début de
	l'histoire, lorsque nous en sommes arrivés
	si loin, les roues tournent à vide et
	nous n'écoutons plus une description de
	quoi que ce soit qui soit le cas ?
	\citep[p.~268]{Boolos2000}
\end{quote}
Si nous avons effectivement dépassé
«~quoi que ce soit qui soit le cas~»,
c'est directement au remplacement qu'il faut
en imputer la responsabilité. C'est en effet
cet axiome qui rend possible notre démonstration.
Dans ce cas, on peut supposer que Boolos
devrait revenir sur son affirmation,
formulée quelques décennies plus tôt,
selon laquelle le remplacement n'a
«~aucune conséquence indésirable~»
(voir la \olref[replacement][extrinsic]{sec}).

Mais l'existence de $\kappa$ est-elle si
inquiétante ? Le \emph{paradoxe de Tristram
Shandy}, dû à Russell, peut nous aider à
y réfléchir. Tristram Shandy tient le journal
de sa vie, mais il lui faut une année pour
raconter une seule journée. Chaque année,
il prend davantage de retard : au bout
d'un an, il n'a raconté qu'un jour et
en a vécu 364 autres dont il n'a rien écrit ;
au bout de deux ans, il n'en a raconté
que deux et en a vécu 728 autres sans
les raconter ; au bout de trois ans,
il n'en a raconté que trois et en a
vécu 1092 autres sans les raconter
\dots\footnote{En négligeant les années
bissextiles.} Pourtant, si Tristram est
\emph{immortel}, il finira par raconter
chaque journée, puisqu'il racontera
la $n$-ième journée pendant la $n$-ième
année de sa vie. Ainsi, «~à la fin
des temps~», son journal sera complet.

Pourquoi cette idée serait-elle si différente
de celle selon laquelle $\alpha$ est plus
petit que $\aleph_\alpha$ — et même de
plus en plus désespérément plus petit —
jusqu'à $\kappa$, où l'écart se résorbe
et où $\kappa = \aleph_\kappa$ ?

Laissons cette question de côté. En supposant
que nous acceptions $\ZFC$, terminons par
quelques autres constructions de points fixes.
Les trois résultats suivants établissent,
intuitivement, qu'il existe un point non trivial
où la hiérarchie est aussi large que haute :

\begin{prop}\ollabel{bethfixed}
Il existe un point fixe de la fonction
$\beth$, c'est-à-dire un $\kappa$ tel que
$\kappa = \beth_\kappa$.
\end{prop}

\begin{proof}
On procède comme dans la \olref{alephfixed},
en remplaçant «~$\aleph$~» par
«~$\beth$~».
\end{proof}

\begin{prop}\ollabel{stagesize}
$\card{V_{\omega+\alpha}} = \beth_{\alpha}$.
Si $\omega \ordtimes \omega \leq \alpha$,
alors $\card{V_\alpha} = \beth_\alpha$.
\end{prop}

\begin{proof}
La première assertion se démontre par une
récurrence transfinie simple. La seconde
en découle, car
$\omega \ordtimes \omega \leq \alpha$
implique $\omega + \alpha = \alpha$.
Pour l'établir, utilisons les propriétés
de l'arithmétique ordinale démontrées dans le
\olref[ord-arithmetic][]{chap}.
Remarquons d'abord que
$\omega \ordtimes \omega
= \omega \ordtimes (1 \ordplus \omega)
= (\omega \ordtimes 1)
\ordplus (\omega\ordtimes\omega)
= \omega \ordplus (\omega \ordtimes \omega)$.
Si $\omega \ordtimes \omega \leq \alpha$,
c'est-à-dire si
$\alpha = (\omega\ordtimes\omega)
\ordplus \beta$ pour un certain $\beta$,
alors
$\omega \ordplus \alpha
= \omega \ordplus
((\omega \ordtimes \omega) \ordplus \beta)
= (\omega \ordplus (\omega
\ordtimes \omega)) \ordplus \beta
= (\omega \ordtimes \omega)
\ordplus \beta = \alpha$.
\end{proof}

\begin{cor}
Il existe un $\kappa$ tel que
$\card{V_\kappa} = \kappa$.
\end{cor}

\begin{proof}
Soit $\kappa$ un point fixe de la fonction
$\beth$, donné par la \olref{bethfixed}.
Il est clair que
$\omega \ordtimes \omega < \kappa$.
Ainsi,
$\card{V_\kappa}
= \beth_\kappa = \kappa$
par la \olref{stagesize}.
\end{proof}

Il y a autant d'étapes avant $V_\kappa$
que d'!!{element}s de $V_\kappa$.
Intuitivement, la hiérarchie jusqu'à
$V_\kappa$ est donc aussi large que haute.
Cela rappelle
l'histoire de Tristram Shandy : nous
passons d'une étape à la suivante
en formant l'\emph{ensemble des parties},
ce qui rend chaque fois la hiérarchie
\emph{beaucoup} plus vaste. Pourtant,
«~à la fin~», c'est-à-dire à l'étape
$\kappa$, sa largeur rattrape sa hauteur.

On peut se demander : \emph{à quelle
fréquence la largeur de la hiérarchie
égale-t-elle sa hauteur ?}
La réponse est : \emph{aussi souvent
qu'il y a d'ordinaux}. Mais cela
demande quelques explications.

Définissons un terme $\tau$ comme suit.
Pour tout $A$, posons :\footnote{Note de l'édition française : l'original initialise $\tau_0(A)=\card{A}$ ; si $\card{A}$ est déjà un point fixe de $\beth$, la suite est constante et $\tau(A)=\card{A}$, de sorte que $W$ ne progresse plus. Le successeur ordinal $\card{A}\ordplus1$ garantit ici $\card{A}<\tau(A)$.}
\begin{align*}
	\tau_0(A) & \defis \card{A}\ordplus 1\\
	\tau_{n+1}(A) & \defis \beth_{\tau_n(A)}\\
	\tau(A) & \defis \bigcup_{n < \omega}\tau_n(A)
\intertext{Comme dans la \olref{bethfixed},
$\tau(A)$ est un point fixe de la fonction
$\beth$ pour tout $A$, et il est clair que
$\card{A} < \tau(A)$. Considérons alors
la définition récursive suivante :}
	W_0 &\defis 0\\
	W_{\alpha + 1} & \defis \tau(W_\alpha)\\
	W_\alpha & \defis \bigcup_{\beta < \alpha} W_\beta
	\text{, si $\alpha$ est un ordinal limite}
\end{align*}
Cette construction est définie pour tous
les ordinaux. Intuitivement, $W$ est donc
une «~!!a{injection}~» des ordinaux
non nuls dans les points fixes de la
fonction $\beth$.\footnote{Note de l'édition
française : l'original inclut implicitement
$W_0=0$ parmi les points fixes de la fonction
$\beth$ ; or $\beth_0=\omega\neq0$.
Seuls les $W_\alpha$ d'indice non nul
appartiennent à cette classe.}
Et, exactement comme auparavant,
la hiérarchie jusqu'à $V_{W_\alpha}$
est aussi large que haute pour tout
$\alpha$ ; le cas $\alpha=0$
est l'égalité triviale $\card{V_0}=0$.

\end{document}
